% Candidate family-format lemma, 2026-11-19, maximal_locus2.
% Not yet independently audited or promoted. Continuity sewing is separate.
\subsection*{Uniform formats for individual surface approximants in a definable family}

Fix an o-minimal expansion $\mathcal R$ of the real field. Let $A$ be a
definable parameter set, let $F,G$ be compact definable closed surfaces,
and let $q:A\times F\to G$ be definable. Assume each $q_a$ is a continuous
cell-like surjection. Joint continuity in $a$ is not required for the
format lemma, although it will be required for the intended sewing problem.

If $A=\varnothing$, the asserted family is empty. If $F=G=\varnothing$,
use the unique empty homeomorphism for every parameter and take its uniform
error to be zero. If exactly one of $F,G$ is empty and $A$ is nonempty, no
surjection $q_a:F\to G$ as in the hypotheses exists. Thus the proof below
may assume $A,F,G$ nonempty.

\paragraph{Candidate lemma.}
There is a single definable map
\[
 \psi:A\times(0,1]\times F\longrightarrow G
 \tag{US.1}
\]
such that, for every $(a,\eta)$, the slice $\psi_{a,\eta}:F\to G$ is a
homeomorphism and
\[
 \sup_{x\in F}\|\psi_{a,\eta}(x)-q_a(x)\|<\eta.
 \tag{US.2}
\]
The map $\psi$ need not be jointly continuous in $a$ or $\eta$. Its graph
is nevertheless defined by one formula with finitely many fixed parameters.
The claim is consequently a genuine finite-format existence assertion,
stronger than independent existence of an approximant for each $(a,\eta)$.

\paragraph{Proof, step 1: first-order scalar-triangulation certificates.}
For each $a$ form the compact definable cylinder $C_a=C_{q_a}$ and its scalar
height $\tau_a:C_a\to[0,1]$ as in A5.016. The predicates that the slices
$q_a$ are continuous surjections and the sets $C_a$ are closed and bounded
are first-order. They therefore persist in elementary extensions of
$\mathcal R$. No cell-likeness assertion is used at this step.

Consider all possible finite certificates for exact scalar triangulations
of these heights. A certificate specifies:
\begin{enumerate}
\item a finite abstract simplicial complex $K$ and coordinates for a geometric
realization $P$ of $K$;
\item a fixed formula $\varphi(u,z;b)$, with a finite tuple $b$ of freely
quantified coefficients, defining the graph of a homeomorphism
$\pi:P\to C_a$;
\item vertex heights, with the assertion that $\tau_a\circ\pi$ is their
affine interpolation on each simplex.
\end{enumerate}
Here is an explicit certificate predicate. Let $K$ have vertices
$1,\ldots,N$ and let $v=(v_1,\ldots,v_N)$ be a tuple of points in
$\mathbb R^D$. The finite integer $D$ is part of the certificate format
$i$, together with $K$ and the graph formula; it is not a quantified field
variable. Scalar triangulation in fact permits the dimension of the ambient
space containing $C_a$, but no fixed-dimension restriction is needed here.
Write
\[
 P_K(v;u)\ \Longleftrightarrow\
 \bigvee_{\sigma\in K}\exists(\lambda_j)_{j\in\sigma}
 \left(\lambda_j\ge0,\ \sum_{j\in\sigma}\lambda_j=1,
       \ u=\sum_{j\in\sigma}\lambda_jv_j\right).
\]
Let $\operatorname{Real}_K(v)$ assert affine independence on each simplex
and that two realized simplices intersect exactly in the realization of
their common face. These are finite first-order conditions; for an empty
common face the intersection condition asserts disjointness. For a fixed
formula $\varphi_i(u,z;b)$, with a fixed finite length for the additional
tuple $b$, define
\[
 \Gamma_i(a,v,b;u,z)
 \ \Longleftrightarrow\
 P_K(v;u)\ \wedge\ z\in C_a\ \wedge\varphi_i(u,z;b).
\]
Let $\operatorname{Homeo}(\Gamma_i)$ express that this relation is a
bijective graph between the indicated sets and that both it and its
inverse are continuous. Continuity is expressed in the usual
$\forall x\,\forall\varepsilon>0\,\exists\delta>0\,\forall x'$
form, using squared Euclidean distances if desired. We can now take
\[
\begin{split}
 C_i(a)\ \Longleftrightarrow\ 
 \exists v\,\exists b\,\exists(c_1,\ldots,c_N)\in[0,1]^N\;
 &\bigl[\operatorname{Real}_K(v)\wedge
        \operatorname{Homeo}(\Gamma_i)\bigr.\\
 &\bigl.\wedge\!\bigwedge_{\sigma\in K}
 \forall\lambda\in\Delta_\sigma\,\forall z\,
 \bigl(\Gamma_i(a,v,b;\textstyle\sum_{j\in\sigma}\lambda_jv_j,z)
 \Rightarrow \tau_a(z)=\textstyle\sum_{j\in\sigma}\lambda_jc_j\bigr)
 \bigr].
\end{split}
\tag{US.2a}
\]
Here $\Delta_\sigma$ is the barycentric simplex defined by nonnegative
coordinates with sum one. In particular, the single-vertex cases force
$c_j=\tau_a(\pi(v_j))$. Every quantifier displayed ranges over a finite
tuple of elements of the field. The index $i$ specifies only the finite
complex, source ambient dimension, and graph formula, not the values of
their coefficients.

Shiota's exact scalar triangulation theorem holds in every elementary
extension of $\mathcal R$, since every such extension is again an
o-minimal expansion of a real closed field. To spell out the latter point,
for each formula in one ordered-field variable, o-minimal cell decomposition
in $\mathcal R$ supplies a uniform finite bound on the number of interval
and point components of its fibres. The sentence giving that bound, with
the other variables universally quantified, belongs to the complete theory
and hence holds in every elementary equivalent model. This applies to
every formula. Therefore each element $a$ of the parameter set in every
such extension satisfies at least one $C_i$.

We claim that finitely many $C_i$ cover $A$ already in $\mathcal R$.
If not, the type
\[
 \{a\in A\}\cup\{\neg C_i(a): i\text{ a certificate format}\}
 \tag{US.3}
\]
would be finitely satisfiable in $\mathcal R$. Compactness, applied with
the elementary diagram of $\mathcal R$, would produce an elementary
extension containing a realization of this type. Scalar triangulation
there supplies a finite complex and a definable homeomorphism with a
finite parameter tuple, hence one of the listed certificate formats,
a contradiction. Thus a fixed finite list of formulas and complexes
suffices. The indexing collection is the set of all finite formulas and
finite complexes. We assume neither that this set is countable nor that
its certificates can be effectively enumerated or decided; logical
compactness applies to this arbitrary set of formulas. This argument
uses the all-model scalar theorem; it does not
infer a format bound from the pointwise real-surface approximation theorem.

\paragraph{Step 2: finitely many abstract pairs and fixed collars.}
Use definable choice in this finite certificate family to select, for every
$a$, a certificate and its coordinates. No continuity of these choices is
asserted. Refine the finite definable partition by the set of vertices
whose chosen scalar height is zero (and, if desired, by those at height
one). There are only finitely many such patterns.

Use the standard abstract realization $|K|$ with basis-vector vertices and
the simplicial coordinate map $\rho_a:|K|\to P_a$. On each part of the
finite partition this is one definable family. Put
\[
 \Pi_a=\pi_a\circ\rho_a,
 \qquad \ell_a=\tau_a\circ\Pi_a.
\]
The domain complex is now fixed; $\ell_a$ is PL on each of its simplices
and varies only through its finitely many vertex heights. Its zero set is
one fixed subcomplex $L\subset K$ on the chosen partition part.

For every real $a$ under consideration, the classical cell-like argument
in A5.016 makes $C_a$ a topological three-manifold with its two extreme
levels as boundary. The finite link proof in A5.016 then says that $K$ is
a PL three-manifold and $L$ a union of boundary components. Consequently
each of the finitely many abstract pairs $(K,L)$ that actually occurs
admits a fixed PL collar
\[
 c_{K,L}:|L|\times[0,1]\longrightarrow|K|.
\]
Choose such a collar once for each pair. This is a finite list of choices,
not a parameter-dependent selection of arbitrarily complicated collars.
We make no assertion about cell-likeness or this topological conclusion in
an arbitrary elementary extension; only the preceding scalar-format step
used such extensions.

\paragraph{Step 3: fixed subdivisions and rational stellar choices.}
For a fixed pair, triangulate the collar once so that its images lie in
simplices of $K$, its source height is linear, and its own graph is
simplicial. Because every $\ell_a$ is linear on the same original $K$,
this one refinement makes all $\ell_a$ linear on the required image
simplices. Its zero-vertex pattern is fixed.

Choose $\epsilon(a)>0$ below all positive source and range vertex heights
required in Lemma 1 of A5.016, and small enough that
$\ell_a^{-1}([0,\epsilon(a)])$ lies in the fixed collar image. This can
be done definably: the complement of the open collar is a fixed compact
polyhedron disjoint from $L$, and the minimum of the PL function $\ell_a$
there is a positive definable function. A common fixed triangulation
of that complement even makes this minimum the minimum of finitely many
affine combinations of the original vertex heights.

The stellar choices in the height-leveling proof can now be made with
fixed rational formulas. For a mixed simplex with zero-face barycenter
$z$ and positive-face barycenter $w$, let $d_a=\ell_a(w)>\epsilon(a)$.
Then
\[
 (1-\epsilon(a)/d_a)z+(\epsilon(a)/d_a)w
 \tag{US.4}
\]
is an interior point of that simplex at the desired height. The analogous
formula uses the fixed source height on the collar domain. The same finite
list of mixed faces is starred in the same order for every $a$ on the
partition part. The resulting canonical simplicial isomorphism is a
definable family: its affine coefficients are obtained by finite linear
algebra from these finitely many moving vertices. All denominators are
nonzero because the prescribed simplices remain nondegenerate.

Thus we obtain a single definable family of height-preserving collars
\[
 C_a:|L|\times[0,\epsilon(a)]
       \longrightarrow\ell_a^{-1}([0,\epsilon(a)]).
\]
For each fixed $a$ it is continuous in both displayed variables. It need
not be continuous as $a$ varies, because the chosen scalar triangulation
$\Pi_a$ need not vary continuously.

\paragraph{Step 4: definable choice of an error scale.}
Repeat the coordinate extraction in A5.016 for this family of collars.
The identification $G\to|L|$ at height zero is definable in $a$ because
it is the inverse of $\Pi_a$ on the zero fibre. We obtain one definable
family of homeomorphisms
\[
 h_{a,s}:F\longrightarrow G,
 \qquad 0<s\le\epsilon(a),
\]
with $h_{a,s}\to q_a$ uniformly on $F$ as $s\downarrow0$ for each fixed
$a$. The error
\[
 E(a,s)=\sup_{x\in F}\|h_{a,s}(x)-q_a(x)\|
\]
is a definable function, since $F$ is compact and the slices are
continuous. The relation
\[
 0<s\le\epsilon(a),\qquad E(a,s)<\eta
\]
has a solution for every $(a,\eta)\in A\times(0,1]$. Definable choice
provides a single function $s(a,\eta)$ satisfying it. Setting
$\psi_{a,\eta}=h_{a,s(a,\eta)}$ proves (US.1)--(US.2).
\hfill$\square$

\paragraph{What this still does not prove.}
Definable choice does not give joint continuity in $(a,\eta)$. For the
family mapping-cylinder application, the separate remaining task is to
replace $\psi$ by a jointly continuous definable family while retaining
uniform convergence to $q_a$ as the error parameter tends to zero,
locally uniformly in $a$.

A useful quantitative identity is available for sewing. If
$u,v:F\to G$ are homeomorphisms, then for $d=uv^{-1}:G\to G$,
\[
 \|d-\operatorname{id}_G\|_\infty
   =\|u-v\|_\infty
   =\|d^{-1}-\operatorname{id}_G\|_\infty.
 \tag{US.5}
\]
Thus approximants close to the same possibly noninjective $q_a$ give
transition homeomorphisms close to the identity; inversion of $q_a$ is
not required. These transition maps lie in a fixed definable family
whenever $u,v$ do. The fixed-format local contraction theorem may apply
to that family, but the finite closure-pool and collar-radius construction
needed to sew all parameter cells has not been proved in this note.

The format lemma is presently a candidate awaiting independent audit.
Neither it nor the individual theorem is assigned original A.5 gate
movement without a proved, appropriately scoped uniform conclusion.
